\documentclass[10pt,a4paper]{amsart}
\usepackage{amsmath,amssymb,amsthm,graphicx}
\usepackage{verbatim}
\newenvironment{reference}{\comment}{\endcomment}
%\newenvironment{reference}{}{}
\newenvironment{slogan}{\comment}{\endcomment}
\newenvironment{history}{\comment}{\endcomment}
% For commutative diagrams we use Xy-pic
\usepackage[all]{xy}

% We use 2cell for 2-commutative diagrams.
\xyoption{2cell}
% Theorem environments.
%
\theoremstyle{plain}
\newtheorem{theorem}[subsection]{Theorem}
\newtheorem{proposition}[subsection]{Proposition}
\newtheorem{lemma}[subsection]{Lemma}

\theoremstyle{definition}
\newtheorem{definition}[subsection]{Definition}
\newtheorem{example}[subsection]{Example}
\newtheorem{exercise}[subsection]{Exercise}
\newtheorem{situation}[subsection]{Situation}

\theoremstyle{remark}
\newtheorem{remark}[subsection]{Remark}
\newtheorem{remarks}[subsection]{Remarks}

\numberwithin{equation}{subsection}

% Macros
%
\def\lim{\mathop{\mathrm{lim}}\nolimits}
\def\colim{\mathop{\mathrm{colim}}\nolimits}
\def\Spec{\mathop{\mathrm{Spec}}}
\def\Hom{\mathop{\mathrm{Hom}}\nolimits}
\def\Ext{\mathop{\mathrm{Ext}}\nolimits}
\def\SheafHom{\mathop{\mathcal{H}\!\mathit{om}}\nolimits}
\def\SheafExt{\mathop{\mathcal{E}\!\mathit{xt}}\nolimits}
\def\Sch{\mathit{Sch}}
\def\Mor{\mathop{\mathrm{Mor}}\nolimits}
\def\Ob{\mathop{\mathrm{Ob}}\nolimits}
\def\Sh{\mathop{\mathit{Sh}}\nolimits}
\def\NL{\mathop{N\!L}\nolimits}
\def\CH{\mathop{\mathrm{CH}}\nolimits}
\def\proetale{{pro\text{-}\acute{e}tale}}
\def\etale{{\acute{e}tale}}
\def\QCoh{\mathit{QCoh}}
\def\Ker{\mathop{\mathrm{Ker}}}
\def\Im{\mathop{\mathrm{Im}}}
\def\Coker{\mathop{\mathrm{Coker}}}
\def\Coim{\mathop{\mathrm{Coim}}}

% Boxtimes
%
\DeclareMathSymbol{\boxtimes}{\mathbin}{AMSa}{"02}

%
% Macros for moduli stacks/spaces
%
\def\QCohstack{\mathcal{QC}\!\mathit{oh}}
\def\Cohstack{\mathcal{C}\!\mathit{oh}}
\def\Spacesstack{\mathcal{S}\!\mathit{paces}}
\def\Quotfunctor{\mathrm{Quot}}
\def\Hilbfunctor{\mathrm{Hilb}}
\def\Curvesstack{\mathcal{C}\!\mathit{urves}}
\def\Polarizedstack{\mathcal{P}\!\mathit{olarized}}
\def\Complexesstack{\mathcal{C}\!\mathit{omplexes}}
% \Pic is the operator that assigns to X its picard group, usage \Pic(X)
% \Picardstack_{X/B} denotes the Picard stack of X over B
% \Picardfunctor_{X/B} denotes the Picard functor of X over B
\def\Pic{\mathop{\mathrm{Pic}}\nolimits}
\def\Picardstack{\mathcal{P}\!\mathit{ic}}
\def\Picardfunctor{\mathrm{Pic}}
\def\Deformationcategory{\mathcal{D}\!\mathit{ef}}
\usepackage{cleveref}
\crefname{originalreference}{original reference}{original references}
\Crefname{originalreference}{Original reference}{Original references}
\makeatletter
\newlabel{divisors-section-blowing-up}{{divisors-section-blowing-up}{source}{}{sourcepointer-1}{}}
\newlabel{divisors-section-blowing-up@cref}{{[originalreference][0][]divisors-section-blowing-up}{[1][0][]source}{}{}{}}
\makeatother
\crefname{originalreference}{original reference}{original references}
\Crefname{originalreference}{Original reference}{Original references}
\makeatletter
\newlabel{divisors-lemma-blowing-up-gives-effective-Cartier-divisor}{{divisors-lemma-blowing-up-gives-effective-Cartier-divisor}{source}{}{sourcepointer-2}{}}
\newlabel{divisors-lemma-blowing-up-gives-effective-Cartier-divisor@cref}{{[originalreference][0][]divisors-lemma-blowing-up-gives-effective-Cartier-divisor}{[1][0][]source}{}{}{}}
\makeatother
\crefname{originalreference}{original reference}{original references}
\Crefname{originalreference}{Original reference}{Original references}
\makeatletter
\newlabel{divisors-lemma-complement-locally-principal-closed-subscheme}{{divisors-lemma-complement-locally-principal-closed-subscheme}{source}{}{sourcepointer-3}{}}
\newlabel{divisors-lemma-complement-locally-principal-closed-subscheme@cref}{{[originalreference][0][]divisors-lemma-complement-locally-principal-closed-subscheme}{[1][0][]source}{}{}{}}
\makeatother
\crefname{originalreference}{original reference}{original references}
\Crefname{originalreference}{Original reference}{Original references}
\makeatletter
\newlabel{coherent-lemma-relative-affine-vanishing}{{coherent-lemma-relative-affine-vanishing}{source}{}{sourcepointer-4}{}}
\newlabel{coherent-lemma-relative-affine-vanishing@cref}{{[originalreference][0][]coherent-lemma-relative-affine-vanishing}{[1][0][]source}{}{}{}}
\makeatother
\crefname{originalreference}{original reference}{original references}
\Crefname{originalreference}{Original reference}{Original references}
\makeatletter
\newlabel{limits-lemma-higher-direct-images-zero-above-dimension-fibre}{{limits-lemma-higher-direct-images-zero-above-dimension-fibre}{source}{}{sourcepointer-5}{}}
\newlabel{limits-lemma-higher-direct-images-zero-above-dimension-fibre@cref}{{[originalreference][0][]limits-lemma-higher-direct-images-zero-above-dimension-fibre}{[1][0][]source}{}{}{}}
\makeatother
\crefname{originalreference}{original reference}{original references}
\Crefname{originalreference}{Original reference}{Original references}
\makeatletter
\newlabel{cohomology-lemma-relative-Leray}{{cohomology-lemma-relative-Leray}{source}{}{sourcepointer-6}{}}
\newlabel{cohomology-lemma-relative-Leray@cref}{{[originalreference][0][]cohomology-lemma-relative-Leray}{[1][0][]source}{}{}{}}
\makeatother
\crefname{originalreference}{original reference}{original references}
\Crefname{originalreference}{Original reference}{Original references}
\makeatletter
\newlabel{cohomology-lemma-apply-Leray}{{cohomology-lemma-apply-Leray}{source}{}{sourcepointer-7}{}}
\newlabel{cohomology-lemma-apply-Leray@cref}{{[originalreference][0][]cohomology-lemma-apply-Leray}{[1][0][]source}{}{}{}}
\makeatother
\title{A blowing up whose fibres have dimension at most d - 1 bounds the cohomological dimension of the ideal by d}
\author{The Stacks Project Authors}
\date{stacks-project:local-cohomology.tex:local-snapshot-20261012; GNU FDL 1.2 or later}
\begin{document}
\maketitle
\noindent\textit{Editorial note.} All delivered mathematical text is the author's own, copied from the pinned Stacks Project chapter source sources/stacks-local-cohomology/pinned-original-source.tex; no statement or proof step was rewritten, reordered or removed.\par\medskip
\noindent\textit{Editorial note.} Preamble provenance: the Stacks Project keeps one single shared preamble for every chapter; this card takes it directly from the staged sources/stacks-local-cohomology/preamble.tex as three verbatim spans (lines 8-12, 14-18 and 154-217), so no derived file was written and nothing was added or redefined. The wrapper class and the base packages are supplied by the assembler.\par\medskip
\noindent\textit{Editorial note.} Bibliography provenance: neither the counted statement nor its proof contains a citation (this lane rescans the whole span before writing the card), so the card carries no bibliography and no entry was reconstructed.\par\medskip
\noindent\textit{Editorial note.} Selection provenance: this card was opened under the file-based control contract (CONTROL.md), with the mechanical pre-card gate run on its statement and proof interval before the card was written.\par\medskip
\noindent\textit{Allowed standard prerequisites (editorial).}\begin{itemize}
\item The complement of an effective Cartier divisor is affine, applied to the exceptional divisor of the blowing up. Invoked inside the counted proof (source lines 879-886) as cross-chapter pointers to the chapter on divisors. Makes the inclusion of the complement of the exceptional divisor affine, so that its higher direct images vanish.
\item The vanishing of the higher direct images of a morphism above the dimension of its fibres, and the Leray spectral sequence. Both are invoked inside the counted proof (source lines 891-898) as cross-chapter pointers. Combine the two vanishing statements into the vanishing of the cohomology of the complement above d, which is the definition of the required bound.
\end{itemize}\medskip
\ifdefined\hypertarget\hypertarget{sourcepointer-1}{}\fi
\noindent\textit{Source reference.} Original reference divisors-section-blowing-up: sources/stacks-local-cohomology/pinned-original-source.tex, source lines 880--880. Navigation pointer only: the counted text refers to the Stacks item divisors-section-blowing-up, which belongs to another chapter of the project and is not part of this card. The pointer carries the target's own label as its number because this chapter's staged snapshot does not contain the other chapter's tag table; its content is not reproduced here.\par\medskip
\ifdefined\hypertarget\hypertarget{sourcepointer-2}{}\fi
\noindent\textit{Source reference.} Original reference divisors-lemma-blowing-up-gives-effective-Cartier-divisor: sources/stacks-local-cohomology/pinned-original-source.tex, source lines 883--883. Navigation pointer only: the counted text refers to the Stacks item divisors-lemma-blowing-up-gives-effective-Cartier-divisor, which belongs to another chapter of the project and is not part of this card. The pointer carries the target's own label as its number because this chapter's staged snapshot does not contain the other chapter's tag table; its content is not reproduced here.\par\medskip
\ifdefined\hypertarget\hypertarget{sourcepointer-3}{}\fi
\noindent\textit{Source reference.} Original reference divisors-lemma-complement-locally-principal-closed-subscheme: sources/stacks-local-cohomology/pinned-original-source.tex, source lines 886--886. Navigation pointer only: the counted text refers to the Stacks item divisors-lemma-complement-locally-principal-closed-subscheme, which belongs to another chapter of the project and is not part of this card. The pointer carries the target's own label as its number because this chapter's staged snapshot does not contain the other chapter's tag table; its content is not reproduced here.\par\medskip
\ifdefined\hypertarget\hypertarget{sourcepointer-4}{}\fi
\noindent\textit{Source reference.} Original reference coherent-lemma-relative-affine-vanishing: sources/stacks-local-cohomology/pinned-original-source.tex, source lines 890--890. Navigation pointer only: the counted text refers to the Stacks item coherent-lemma-relative-affine-vanishing, which belongs to another chapter of the project and is not part of this card. The pointer carries the target's own label as its number because this chapter's staged snapshot does not contain the other chapter's tag table; its content is not reproduced here.\par\medskip
\ifdefined\hypertarget\hypertarget{sourcepointer-5}{}\fi
\noindent\textit{Source reference.} Original reference limits-lemma-higher-direct-images-zero-above-dimension-fibre: sources/stacks-local-cohomology/pinned-original-source.tex, source lines 893--893. Navigation pointer only: the counted text refers to the Stacks item limits-lemma-higher-direct-images-zero-above-dimension-fibre, which belongs to another chapter of the project and is not part of this card. The pointer carries the target's own label as its number because this chapter's staged snapshot does not contain the other chapter's tag table; its content is not reproduced here.\par\medskip
\ifdefined\hypertarget\hypertarget{sourcepointer-6}{}\fi
\noindent\textit{Source reference.} Original reference cohomology-lemma-relative-Leray: sources/stacks-local-cohomology/pinned-original-source.tex, source lines 895--895. Navigation pointer only: the counted text refers to the Stacks item cohomology-lemma-relative-Leray, which belongs to another chapter of the project and is not part of this card. The pointer carries the target's own label as its number because this chapter's staged snapshot does not contain the other chapter's tag table; its content is not reproduced here.\par\medskip
\ifdefined\hypertarget\hypertarget{sourcepointer-7}{}\fi
\noindent\textit{Source reference.} Original reference cohomology-lemma-apply-Leray: sources/stacks-local-cohomology/pinned-original-source.tex, source lines 898--898. Navigation pointer only: the counted text refers to the Stacks item cohomology-lemma-apply-Leray, which belongs to another chapter of the project and is not part of this card. The pointer carries the target's own label as its number because this chapter's staged snapshot does not contain the other chapter's tag table; its content is not reproduced here.\par\medskip
\begin{lemma}
\label{lemma-cd-blowup}
Let $A$ be a Noetherian ring. Let $I \subset A$ be an ideal.
Let $b : X' \to X = \Spec(A)$ be the blowing up of $I$.
If the fibres of $b$ have dimension $\leq d - 1$, then
$\text{cd}(A, I) \leq d$.
\end{lemma}

\begin{proof}
Set $U = X \setminus V(I)$. Denote $j : U \to X'$ the canonical open
immersion, see Divisors, Section \ref{divisors-section-blowing-up}.
Since the exceptional divisor is an effective Cartier divisor
(Divisors, Lemma
\ref{divisors-lemma-blowing-up-gives-effective-Cartier-divisor})
we see that $j$ is affine, see
Divisors, Lemma
\ref{divisors-lemma-complement-locally-principal-closed-subscheme}.
Let $\mathcal{F}$ be a quasi-coherent $\mathcal{O}_U$-module.
Then $R^pj_*\mathcal{F} = 0$ for $p > 0$, see
Cohomology of Schemes, Lemma
\ref{coherent-lemma-relative-affine-vanishing}.
On the other hand, we have $R^qb_*(j_*\mathcal{F}) = 0$ for
$q \geq d$ by Limits, Lemma
\ref{limits-lemma-higher-direct-images-zero-above-dimension-fibre}.
Thus by the Leray spectral sequence
(Cohomology, Lemma \ref{cohomology-lemma-relative-Leray})
we conclude that $R^n(b \circ j)_*\mathcal{F} = 0$ for
$n \geq d$. Thus $H^n(U, \mathcal{F}) = 0$ for $n \geq d$
(by Cohomology, Lemma \ref{cohomology-lemma-apply-Leray}).
This means that $\text{cd}(A, I) \leq d$ by definition.
\end{proof}

\end{document}
